% =============================================================================
% PREFACE — Deep Reinforcement Learning
% =============================================================================
\chapter*{Preface}
\addcontentsline{toc}{chapter}{Preface}
\markboth{Preface}{Preface}

\section*{Objectives of This Book}

Deep Reinforcement Learning (DRL) combines classical reinforcement learning with
deep learning. Over the past decade, it has enabled spectacular advances:
defeating human champions at Go and Atari games, controlling robotic manipulators
with remarkable dexterity, optimizing financial strategies, and designing new
therapeutic molecules.

This book is intended for Master's and PhD students in computer science, applied
mathematics, or engineering, as well as researchers and practitioners seeking a
rigorous understanding of the theoretical foundations and algorithmic methods of DRL.

\section*{Course Organization}

The course is structured in eleven chapters, organized in a pedagogical progression
from mathematical foundations to recent applications:

\begin{enumerate}
    \item \textbf{Markov Decision Processes (MDPs)}: the fundamental mathematical
    framework that formalizes agent-environment interaction.

    \item \textbf{Dynamic Programming}: value iteration and policy iteration
    algorithms, the foundation of all RL methods.

    \item \textbf{Monte Carlo and Temporal Difference Methods}: model-free
    approaches based on sampling.

    \item \textbf{Q-Learning and SARSA}: fundamental tabular algorithms for
    off-policy and on-policy learning.

    \item \textbf{Deep Q-Networks (DQN)}: the revolution of function approximation
    via neural networks applied to RL.

    \item \textbf{Policy Gradient --- REINFORCE}: direct optimization of the
    policy via gradient ascent.

    \item \textbf{Actor-Critic Methods (A3C, PPO)}: combining value estimation
    and policy optimization.

    \item \textbf{State of the Art (SAC, TD3, DDPG)}: modern algorithms for
    continuous action spaces.

    \item \textbf{Multi-Agent Reinforcement Learning}: generalization to
    multi-agent systems.

    \item \textbf{Constrained and Safe RL}: incorporating safety constraints
    into learning.

    \item \textbf{Applications}: robotics, games, finance, and beyond.
\end{enumerate}

\section*{Prerequisites}

The reader should possess solid knowledge in:

\begin{itemize}
    \item \textbf{Probability and Statistics}: random variables, conditional
    expectation, Markov chains, stochastic convergence.

    \item \textbf{Optimization}: gradient descent, convexity, Lagrange multipliers,
    numerical methods.

    \item \textbf{Linear Algebra}: vector spaces, matrix decompositions, eigenvalues.

    \item \textbf{Machine Learning}: neural networks, backpropagation,
    regularization, deep architectures (CNN, RNN).

    \item \textbf{Python Programming}: proficiency with NumPy, PyTorch, and
    Gymnasium environments (formerly OpenAI Gym).
\end{itemize}

\section*{Notation}

We adopt the following conventions throughout this book:

\begin{center}
\renewcommand{\arraystretch}{1.3}
\begin{tabular}{cl}
    \toprule
    \textbf{Symbol} & \textbf{Meaning} \\
    \midrule
    $\mathcal{S}$ & State space \\
    $\mathcal{A}$ & Action space \\
    $\mathcal{R}$ & Reward space \\
    $\gamma \in [0,1)$ & Discount factor \\
    $\pi(a \mid s)$ & Stochastic policy \\
    $V^\pi(s)$ & State-value function under $\pi$ \\
    $Q^\pi(s,a)$ & Action-value function under $\pi$ \\
    $A^\pi(s,a)$ & Advantage function under $\pi$ \\
    $P(s' \mid s, a)$ & Transition probability \\
    $R(s, a, s')$ & Reward function \\
    $\E[\cdot]$ & Expectation \\
    $\PP(\cdot)$ & Probability \\
    $\nabla_\theta$ & Gradient with respect to $\theta$ \\
    $\theta$ & Policy or network parameters \\
    $\alpha$ & Learning rate \\
    $\epsilon$ & Exploration parameter ($\epsilon$-greedy) \\
    $\tau$ & Soft update coefficient \\
    $\mathcal{D}$ & Replay buffer \\
    \bottomrule
\end{tabular}
\end{center}

\section*{Typographic Conventions}

\begin{itemize}
    \item \textbf{Definitions} are presented in blue-bordered boxes.
    \item \textbf{Theorems} and \textbf{propositions} appear in formal boxes with proofs.
    \item \textbf{Algorithms} are presented in pseudocode in dedicated boxes.
    \item \textbf{Implementations} in PyTorch/Gymnasium are given in syntax-highlighted
    code blocks.
    \item \textbf{Key formulas} are highlighted in special boxes.
    \item Points requiring \textbf{special attention} are marked in warning boxes.
\end{itemize}

\section*{Acknowledgments}

This book is the product of several years of teaching and research in reinforcement
learning. We thank the students who, through their questions and curiosity,
contributed to improving the presentation of these topics. We also express our
gratitude to the open-source community, in particular the developers of PyTorch,
Gymnasium, Stable-Baselines3, and CleanRL, whose tools make reinforcement
learning accessible to all.

\section*{How to Use This Book}

\begin{intuition}
Each chapter follows a consistent structure:
\begin{enumerate}
    \item \textbf{Motivation and Intuition} --- Why this method? What problem does
    it solve?
    \item \textbf{Theoretical Foundations} --- Definitions, theorems, convergence proofs.
    \item \textbf{Algorithms} --- Detailed pseudocode with complexity analysis.
    \item \textbf{Implementation} --- Working, reproducible PyTorch code.
    \item \textbf{Exercises} --- Theoretical and practical problems of increasing
    difficulty.
\end{enumerate}
\end{intuition}

\section*{Brief History}

Reinforcement learning has its roots in Bellman's work on dynamic programming
(1950s) and behavioral psychology (Skinner's operant conditioning). Major
milestones include:

\begin{itemize}
    \item \textbf{1989}: Watkins introduces Q-Learning, the first convergent
    off-policy algorithm for MDPs.
    \item \textbf{1992}: Tesauro's TD-Gammon learns backgammon through self-play
    with temporal differences.
    \item \textbf{2013}: Mnih et al.\ (DeepMind) combine DQN with convolutional
    networks to play Atari games from pixels.
    \item \textbf{2015}: DQN published in \emph{Nature}; Silver et al.\ develop
    AlphaGo.
    \item \textbf{2016}: AlphaGo defeats Lee Sedol, world Go champion.
    \item \textbf{2017}: Schulman et al.\ introduce PPO; Haarnoja et al.\ propose SAC.
    \item \textbf{2019}: OpenAI Five defeats world champions at Dota~2.
    \item \textbf{2020--2024}: RL applied to drug discovery, nuclear fusion control
    (DeepMind/EPFL), and LLM alignment (RLHF).
\end{itemize}

\section*{Supplementary Resources}

To deepen the topics covered in this book, we recommend:

\begin{itemize}
    \item \textbf{Sutton \& Barto} --- \emph{Reinforcement Learning: An Introduction}
    (2nd edition, 2018). The standard reference in classical RL.

    \item \textbf{Bertsekas} --- \emph{Dynamic Programming and Optimal Control}
    (4th edition). Rigorous treatment of dynamic programming.

    \item \textbf{Szepesv\'ari} --- \emph{Algorithms for Reinforcement Learning}.
    Concise synthesis of fundamental algorithms.

    \item \textbf{Agarwal et al.} --- \emph{Reinforcement Learning: Theory and
    Algorithms} (2022). Modern theoretical treatment.

    \item \textbf{Online courses}: David Silver (UCL), Sergey Levine (UC Berkeley),
    Emma Brunskill (Stanford).
\end{itemize}

\vfill

\begin{flushright}
\textit{[Author Name]} \\
\textit{March 2026}
\end{flushright}
